\documentclass[11pt,letterpaper]{article}
\usepackage[T1]{fontenc}
\usepackage{lmodern,amsmath,amssymb,booktabs,microtype,xcolor,fancyhdr}
\usepackage[margin=0.90in,top=0.82in,bottom=0.85in,headheight=15pt]{geometry}
\usepackage{titlesec,needspace}
\definecolor{ink}{HTML}{173449}
\definecolor{link}{HTML}{00675E}
\usepackage[colorlinks=true,urlcolor=link,linkcolor=link,pdftitle={What Dirichlet GRH settles beyond the zero theorem}]{hyperref}
\setlength{\parindent}{0pt}
\setlength{\parskip}{4.5pt plus 1pt minus 1pt}
\setlength{\emergencystretch}{2em}
\setcounter{secnumdepth}{1}
\titleformat{\section}{\large\bfseries\color{ink}}{\thesection}{0.65em}{}
\titlespacing{\section}{0pt}{2ex}{1ex}
\titleformat{\paragraph}[runin]{\normalfont\bfseries\color{ink}}{}{0pt}{}
\titlespacing{\paragraph}{0pt}{1ex}{0.45em}
\pagestyle{fancy}\fancyhf{}\lhead{\small\sffamily DIRICHLET GRH}\rhead{\small\sffamily Arithmetic applications}\cfoot{\small\thepage}
\renewcommand{\headrulewidth}{0.35pt}
\raggedbottom
\begin{document}
\thispagestyle{plain}
{\sffamily\small\bfseries\color{link} MATHEMATICAL RESEARCH NOTE\hfill 25 NOVEMBER 2026}\par
\vspace{6pt}
{\fontsize{24}{27}\selectfont\bfseries\color{ink} What Dirichlet GRH settles\\beyond the zero theorem\par}
\vspace{4pt}
{\large Prime cubes, primes in geometric intervals, and isogenies over sextic fields\par}
\medskip
The Lean development of the Alp\"oge--Claude proof in Luke Sernau's
\texttt{grh-lean} repository proves the Riemann hypothesis for ordinary
Dirichlet $L$-functions. This note uses its public theorem to discharge
analytic hypotheses and its one-sided zero-detection argument to isolate a
remaining nonabelian problem. The arithmetic deductions have written proofs
and executable certificates. The verified Lean declarations belong to the
upstream development.

\paragraph{An exact additive threshold.}
Every integer greater than $10{,}500$ is a prime plus a positive multiple
of a prime cube. Exactly $62$ positive integers are exceptions. This
strengthens the prime-square question treated by Lee and O'Clarey.

\paragraph{A published conjecture on geometric intervals.}
For every modulus $q\ge2$, every reduced class occurs among the primes
between $q^t$ and $q^{t+1}$ for every integer $t\ge2$. This completes that
range of Barthel and M\"uller's conjecture. The first interval is also
verified through $q=50{,}000$. A related argument identifies every possible
failure of Pomerance's divisibility conjecture with a least prime in
$1\bmod q$, and gives a GRH-independent finite verification through
$2\cdot10^{14}+1$.

\paragraph{A partial answer to even-degree isogeny uniformity.}
For two families of sextic fields, all odd Momose type 2 isogeny primes lie
in an explicit $66$-prime set with maximum $253{,}507$. In each family a
single field realizes all $66$ as possibilities for the local necessary
condition. The primitive family has normal closure group $A_5$.
A classification using binary expansions extends the uniform bound to
certain alternating-group fields of arbitrarily large even degree.

\paragraph{A boundary for the analytic method.}
For any number field we compute exactly how well a congruence test can predict
complete splitting. For the primitive sextic fields a one-sided inequality
with four explicit prime weights is equivalent to the missing Dedekind-zeta
Riemann hypothesis. The inequality remains to be proved.

\paragraph{Relation to the earlier reports.}
The five \href{https://github.com/VladimirReshetnikov/ProveIt/tree/a92443f00663b1b8b3c233d17df0ad49dc4f2b2c/docs/incoming}{ProveIt incoming archives}
treat CM coefficient fields, an abelian zeta-function classification, quadratic
class groups, primitive roots, norm-Euclidean fields, weighted prime races and auxiliary twists.
Our isogeny question concerns Galois-stable lines on elliptic curves, a
different uniformity problem from classifying CM coefficient fields. The
last section extends the earlier congruence and one-sided tests to arbitrary
number fields. The comparisons establish what is added to those reports.
No claim of publication priority is required by the proofs.
\clearpage
\section{The theorem available for transfer}
\label{sec:lean-audit}

Luke Sernau's repository rewrites the proof attributed to Levent Alp\"oge with Claude.
The readable and compact editions have identical public statement files. Each imports
Mathlib alone. With $q$ a positive integer and $\chi$ any Dirichlet character modulo $q$,
the exported nonvanishing theorem gives
\[
 \Re(s)>\tfrac12,\quad (\chi\ne 1\ \text{or}\ s\ne1)
 \quad\Longrightarrow\quad L(s,\chi)\ne0 .
\]
Here $1$ denotes the principal character. The other four declarations give
Mathlib's Riemann hypothesis, the corresponding explicit statement for $\zeta$,
the critical-strip assertion for every Dirichlet character, and its specialization
to primitive characters. These statements make no hypothesis about central
nonvanishing. The half-plane and strip forms include imprimitive characters.
See the pinned
\href{https://github.com/lsernau/grh-lean/blob/72fead102f663d5c7a407964a13dc0cf9d977207/grh/src/Statements.lean}{public types}
and \href{https://github.com/lsernau/grh-lean/blob/72fead102f663d5c7a407964a13dc0cf9d977207/grh/src/Main.lean}{final declarations}.

\paragraph{How the proof reaches this statement.}
For a multiple $q_0$ of $18480$ and prime $d\equiv q_0/2-1\pmod {q_0}$ the development
compares two estimates of the same lattice covolume with an index correction.
Short integral forms and a spanning estimate give the upper estimate.
A determinant bound for chains of theta lifts gives the lower estimate.
Cancellation of their common main term bounds a weighted prime sum.
After comparison with the prime number theorem at a second cutoff, the estimate becomes
$R_{q_0,d}(x)\le x^{-1/2+\varepsilon}$ for sufficiently large $x$,
where $R_{q_0,d}$ is obtained by centering
\[
  \sum_{\substack{n\le x\\n\equiv1\ (q_0)}}\!
     \Lambda(n)\bigl(1-(n/d)\bigr)\bigl(n^{-1}-x^{-1}\bigr).
\]
For a putative zero $\rho$ in the open right half of the critical strip,
choose $1/2<\beta<\Re\rho$.
The twist-selection result supplies $d$ such that the twisted factors have no
zeros in $\Re s\ge\beta$, $|\Im s|\le|\Im\rho|$. In this window the
numerator of a residue of the transform of $R_{q_0,d}$ counts untwisted zeros
with nonnegative multiplicities. Positivity at real poles and the one-sided bound
allow Landau's lemma to force the residue at $\rho$ to vanish. Finally one
induces an arbitrary character of modulus $q$ to $q_0=18480q$.
An extension of this method to another family would need both the requisite
arithmetic inequality and control of zeros of its auxiliary twists.
See the
\href{https://github.com/lsernau/grh-lean/blob/72fead102f663d5c7a407964a13dc0cf9d977207/grh/src/New/Reduction/MainInequality.lean}{covolume comparison},
\href{https://github.com/lsernau/grh-lean/blob/72fead102f663d5c7a407964a13dc0cf9d977207/grh/src/New/Reduction/OneSidedBound.lean}{remainder theorem}, and
\href{https://github.com/lsernau/grh-lean/blob/72fead102f663d5c7a407964a13dc0cf9d977207/grh/src/New/Reduction/Zeros/Criterion.lean}{zero criterion}.

\paragraph{Verification record.}
The source revision is \texttt{72fead102f663d5c7a407964a13dc0cf9d977207},
using Lean 4.35.0 and Mathlib \texttt{56a00585}. Our SHA-256 check of the
uncompressed \texttt{git archive} agrees with the digest in Andrew Sutherland's
verification record. His published checks cover all 803 readable and 794 compact
source modules. He reports that both builds, per-module kernel replays and the
three independent export checkers nanoda, con-leche and con-ron passed,
and that all five declarations use only \texttt{propext}, \texttt{Classical.choice}
and \texttt{Quot.sound}. The compact edition also passed the full paranoid
Comparator on two machines. The local work for this report consisted of source
inspection and a hash check. The kernel results are Sutherland's. See
\href{https://math.mit.edu/~drew/shortgrh/}{the verification logs and exports}.

\section{A sharp additive threshold for prime cubes}
Lee and O'Clarey's Theorem 1.3 states that Dirichlet GRH implies that
every integer $n>24$ is the sum of a prime and a positive integer
divisible by a prime square. The Lean result supplies their stated
hypothesis. What happens if the second summand must be divisible by a
prime \emph{cube}? An explicit least-prime bound gives a sharp
threshold, and a complete exception list.\hfill
\href{https://doi.org/10.1007/s00013-026-02275-6}{[Lee--O'Clarey, 2026]}

\paragraph{Theorem.} Every integer $n>10500$ has a representation
$n=p+kq^3$ with $p,q$ prime and $k\ge1$. Precisely 62 positive integers
have no such representation. They are
\[
\begin{gathered}
1,2,3,4,5,6,7,8,9,12,14,16,17,20,22,24,28,36,48,52,54,60,62,\\
72,76,78,96,102,108,120,126,150,174,180,216,240,270,294,312,342,\\
420,468,486,528,552,720,750,780,828,870,924,966,1080,1236,1260,\\
1500,1590,1800,1860,2520,8190,10500 .
\end{gathered}
\]
For example, $10501=5+1312\cdot 2^3$. The quantifier on $k$ matters:
requiring the second summand itself to be a cube gives a different problem.

\paragraph{Proof.}
Let $q$ be the least prime which does not divide $n$, and
$B_q=\prod_{\ell<q,\ \ell\text{ prime}}\ell$. Then $B_q\mid n$ and
$n$ is a unit modulo $q^3$. Write $P(a,m)$ for the least prime
congruent to $a\pmod m$. If $P(n,q^3)<n$, subtraction supplies the
required positive multiple. Lamzouri--Li--Soundararajan's Corollary 1.2
applies to every modulus $q^3$ and gives
\[
 P(n,q^3)\le [\varphi(q^3)\log(q^3)]^2<9q^8 .
\]
The estimate uses the Dirichlet GRH supplied by the repository.
\hfill\href{https://arxiv.org/pdf/1309.3595}{[Lamzouri--Li--Soundararajan]}

For every prime $q\ge43$, $B_q>9q^8$. Exact multiplication verifies
this for $43\le q\le257$. If the inequality holds for a prime $q\ge257$
and $r$ is the next prime, then $r<2q$ by Bertrand's postulate and
$B_r=qB_q>9q^9>9(2q)^8>9r^8$. This proves the claim by induction.
Thus every $n$ whose least missing prime is at least 43 has a representation.

For $q\le41$ set $M_q=\max_{(a,q)=1}P(a,q^3)$, with $a$ running
through all residue classes modulo $q^3$. A sieve produces the following
exact values by retaining the first prime in each reduced class.
\begin{center}
\begin{tabular}{r|rrrrrrr}
$q$ &2&3&5&7&11&13&17\\
$M_q$&17&109&1621&9371&47737&120401&349529\\ \hline
$q$&19&23&29&31&37&41&\\
$M_q$&540961&1494349&2367523&3812923&5897123&7802437&
\end{tabular}
\end{center}
An exception must satisfy $B_q\le n\le M_q$. For each of
$q=23,29,31,37,41$ the lower endpoint exceeds the upper. Thus it
suffices to decide representability for $1\le n\le540961$.
For every prime $q$ with $q^3+2\le540961$ and every residue $a\bmod q^3$,
find the first prime $r$ in that class with $r+q^3\le540961$, if one
exists. Mark $r+q^3,r+2q^3,\ldots$ up to the bound. The marked values
have explicit representations. Conversely, every representation occurs
among the marks, since its prime is at least the first prime in the
same class. This step includes the class of $q$ itself. The 62 unmarked
values are those displayed above. The accompanying Python program carries
out the sieve, the primorial checks and the final enumeration. It also
tests each listed exception a second time by subtracting every eligible
positive multiple of each prime cube. An independent bit-set sum of all primes
and all positive cube multiples agrees on every integer through $540961$.

The largest exception can also be checked by hand. For $n=10500$, the
eligible primes $q$ are $2,3,5,7,11,13,17,19$. The first four divide $n$,
so the only candidate prime summand would be $q$ itself, and in each case
$q^3\nmid n-q$. For the remaining four primes an even multiplier would
give an even summand, and $n-2=10498$ is divisible by none of their cubes.
For odd multipliers the possible positive differences, grouped by $q$,
are
\[
 \{9169,6507,3845,1183\},\quad \{8303,3909\},\quad
 \{5587\},\quad \{3641\}.
\]
They have prime divisors $53,3,5,7$, then $19,3$, then $37$ and $11$,
respectively. This gives a direct check of sharpness.

\paragraph{Recovering the square theorem.}
A multiple of $q^3$ is a multiple of $q^2$. Thus only the 62 cube
exceptions can fail to be a prime plus a positive multiple of a prime
square. Direct subtraction of prime-square multiples on these 62 inputs
gives exactly $1,2,3,4,5,8,24$ as square exceptions. This gives a second
proof of Lee and O'Clarey's explicit conclusion, using the same
least-prime estimate as the cube theorem.

For each fixed exponent $j\ge2$ the same argument reduces
$n=p+kq^j$ to finitely many checks. Indeed the least-prime bound is
at most $j^2 q^{2j+2}$, while $B_q$ eventually exceeds every fixed
power of $q$ (for example by the prime number theorem). This observation
does not require a claim of priority for qualitative finiteness. The
cube calculation identifies the last exception exactly.

\section{Primes in geometric intervals}
Barthel and M\"uller ask whether every interval $[q^t,q^{t+1})$ meets every
reduced residue class modulo $q$, for integers $q\ge2$ and $t\ge1$.
Their Lemmas 5--6 cover all $q\le45000$ (with Bertrand's postulate for $q=2$).
Their Remark 7 obtains all $t\ge3$ on Dirichlet GRH, and $t=2$ for
$q\ge17386763$. The following closes the entire intervening range. The
finite computation also extends their $t=1$ verification.
\hfill\href{https://doi.org/10.1080/00029890.2022.2116250}{[Barthel--M\"uller, 2022]}

\paragraph{Theorem.} For every $q\ge2$, every $a$ coprime to $q$, and
every integer $t\ge2$, some prime $p\equiv a\pmod q$ satisfies
$q^t<p<q^{t+1}$. The same conclusion holds for $t=1$ whenever $q\le50000$.

\paragraph{Proof.}
Dudek, Greni\'e and Molteni, Theorem 1.1 and Table 1, prove under Dirichlet
GRH that there is such a prime within distance
\[
 H(x)=q\sqrt x\bigl(0.6265\log x+\log q+14\bigr)
\]
of $x$ if $q\ge3$ and $x\ge(18q\log q)^2$. We have enlarged both their
radius and their threshold using $\varphi(q)\le q$. Their explicit formula
applies to primitive Dirichlet functions and restores the factors at primes
dividing the modulus. This is precisely the zero hypothesis exported by the
Lean repository.\hfill\href{https://doi.org/10.1142/S1793042119500441}{[Dudek--Greni\'e--Molteni, 2019]}

For $q\ge50000$ write $L=\log q$ and $x_2=(18qL)^2$.
The quotient $q/(378L^2+252L)$ is increasing for $L>2$ by
differentiation. At $q=50000$, $10<L<11$, giving
$378L^2+252L<48510<50000$. Thus $x_2<q^3$ and
\[
 \frac{H(x_2)}{q^2}<18L(3L+14),\qquad
 q^2<x_2-H(x_2)<x_2+H(x_2)<q^3.
\]
Indeed the upper endpoint divided by $q^2$ is at most $378L^2+252L$,
and the lower one exceeds $270L^2-252L>1$. For $t\ge2$ set
$x_t=q^{t-2}x_2$. The function $H(x)/x$ decreases for $x\ge x_2$,
again by differentiation. Therefore the interval of radius $H(x_t)$
about $x_t$ is contained in the dilation by $q^{t-2}$ of the interval
just proved to lie inside $(q^2,q^3)$. This proves every $t\ge2$ in the
large-modulus range.

For $45001\le q\le50000$ the included exact sieve verifies $t=1,2$
in every reduced class. The following numbers are maxima of the
\emph{first} prime witnesses when the search starts at the indicated
left endpoint. The endpoints were chosen to let every modulus share a
single sieve.
\begin{center}
\begin{tabular}{rrrr}
$t$ & Search starts at & Classes checked & Largest witness\\ \hline
1 & $50001$ & $144390310$ & $9819037$\\
2 & $2500000001$ & $144390310$ & $2515292327$
\end{tabular}
\end{center}
For each row the program sieves $25000000$ consecutive integers by all
primes through the square root of the last integer. It verifies that the
segment lies inside the appropriate geometric interval for each $q$,
then traverses each reduced progression until it finds a prime. This
uses only integer arithmetic. For these moduli $t\ge3$ follows from
Barthel--M\"uller's Remark 7. Their small-modulus results complete the proof.
Neither endpoint can be a prime in the indicated class for $q\ge2$.

\section{An exact obstruction to Pomerance's divisibility conjecture}
For a prime $p\ge3$, put
\[
 B(p)=\prod_{r<p,\ r\ {\rm prime}}r(r-1),\qquad
 D(p)=\frac{p-1}{\gcd(p-1,B(p))}.
\]
Pomerance conjectured $D(p)=1$ for every $p$. He observed that if the
least prime $P(q)\equiv1\pmod q$ satisfies $q^2\mid P(q)-1$ for a
prime $q$, then $p=P(q)$ violates the conjecture. The converse also
holds when we use the repository's Dirichlet GRH. The proof gives the
exact size of every possible divisibility defect.
\hfill\href{https://math.dartmouth.edu/~carlp/PDF/carmichaelconjecture.pdf}{[Pomerance, 1974]}

\paragraph{Theorem.} Either $D(p)=1$, or $D(p)$ equals a single
prime $q$ for which $p=P(q)$ and $v_q(p-1)=2$. In particular the
conjecture is equivalent to the assertion that
$q^2\nmid P(q)-1$ for every prime $q$.

\paragraph{Proof.}
Fix a prime $q\mid p-1$ and put $e=v_q(p-1)$. The factor $q$ itself
occurs in $B(p)$. If $q\le50000$, the geometric-interval theorem for
$t=1,\ldots,e-1$ supplies $e-1$ distinct primes $r<p$ with
$q\mid r-1$. This handles all $e$, also when the list is empty.

Suppose $q>50000$. Above the class $1\pmod q$ there are two reduced
classes modulo $3q$, by the Chinese remainder theorem. The explicit
least-prime theorem of Lamzouri, Li and Soundararajan (Corollary 1.2)
gives a prime in each such class at most
$[2(q-1)\log(3q)]^2<q^3$. The final inequality follows from
$4\log^2(3q)<q$, valid at $50000$ and thereafter by monotonicity of
$q/\log^2(3q)$. These primes are distinct. When $e\ge3$, use these two
primes first and then one prime from each of the $e-3$ intervals
$[q^j,q^{j+1})$, $j=3,\ldots,e-1$, always in class $1\pmod q$.
They all lie below $q^e\le p-1$. Together with the factor $q$ they give
$v_q(B(p))\ge e$.\hfill\href{https://arxiv.org/pdf/1309.3595}{[Lamzouri--Li--Soundararajan, Cor.~1.2]}

The only possible shortage is $e=2$. It occurs precisely when there
is no earlier prime $r\equiv1\pmod q$, i.e., when $p=P(q)$, and the
shortage is then exactly one. If $q<s$ were two deficient primes at
$p$, we would have
$q^2s^2\le p-1<p=P(q)\le[(q-1)\log q]^2$ by the same
least-prime theorem. This forces $s<\log q<q$, a contradiction.

\paragraph{Frequency and a finite exclusion.}
Every exceptional $p\le X$ has a distinct associated $q$, because $p$
is determined by $q$. Both $p$ and $q$ are odd and $q^2\mid p-1$,
so $2q^2\le p-1$. Thus there are at most
$\pi(\sqrt{(X-1)/2})$ exceptions, in particular
$O(\sqrt X/\log X)$. The relative density of exceptions among primes
is zero. An independent finite check in the archive finds $P(q)<q^2$
for all $664579$ prime moduli $q\le10^7$. It sieves the moduli and
tests candidates $kq+1$ in increasing order by exact trial division.
The largest least prime encountered is $1753933943$, and the largest
multiplier $k$ is $234$, first attained at $q=4722079$. In fact the finite bound
\[
 p\le 2\cdot10^{14}+1\qquad\Longrightarrow\qquad D(p)=1
\]
can be certified independently of GRH. For every prime $q\le10^7$ and every
$e\ge2$ with $q^e\le2\cdot10^{14}$, the same program verifies that the
$(e-1)$st prime in $1\bmod q$ lies below $q^e$. This checks $671350$
prime powers. If $q$ is odd and $q^2\mid p-1$, then $2q^2\le p-1$,
so these ranges include all relevant $q,e$ for the asserted bound, also
for $q=2$. The factor $q$ and the $e-1$ earlier primes supply $q^e$
in $B(p)$. A counterexample would require $q>10^7$ and
$p=1+2kq^2\le[(q-1)\log q]^2$ for some positive integer $k$,
with every smaller candidate prime in $1\pmod q$ absent.

The universally quantified $t=1$ question, the existence of such a $q$,
and Carmichael's conjecture about multiplicities of totients all remain
unresolved by these arguments. Pomerance's divisibility assertion would
rule out one specific proposed construction of a totient with a single
preimage. It would not by itself settle the multiplicity conjecture.

\section{A uniform isogeny bound for two families of sextic fields}
\label{sec:sextic}

Let $E$ be an elliptic curve over a number field $K$. A cyclic isogeny of
prime degree $p$ over $K$ gives a Galois-stable line in $E[p]$. Write
$\lambda$ for the character describing the action on the line, and $\chi_p$
for the mod-$p$ cyclotomic character. We use \emph{Momose type 2} in the
sense of Banwait and Derickx, with $\lambda^{12}=\chi_p^6$ and
$p\equiv3\pmod4$ as in their type 2 sieve. The bounds below concern this
particular type. A bound independent of the discriminant of $K$ is the
relevant uniformity question.

\paragraph{Sextic-field theorem.}
Suppose the normal closure $L$ of $K/\mathbb Q$ acts on the six embeddings
of $K$ as either
$A_4$ on the six edges of a tetrahedron or $A_5$ on its six Sylow
$5$-subgroups. Then every odd Momose type 2 isogeny prime belongs to an
explicit set of $66$ primes, with largest member $253{,}507$. The set
is exactly the set compatible with the local obstruction below for a single field in each family simultaneously, among primes $p\equiv3\pmod4$. Whether the largest prime occurs as
the degree of an actual type 2 isogeny over one of these fields is a separate
existence question. The $A_5$ family consists of primitive sextic fields,
since the stabilizer of a Sylow $5$-subgroup is maximal in $A_5$ (the action
is doubly transitive, as the stabilizer's $5$-cycle permutes the other
five points).

\paragraph{Two inputs.}
Banwait and Derickx, \href{https://doi.org/10.2140/ant.2025.19.1147}
{\emph{Explicit isogenies of prime degree over number fields}},
Proposition 6.1, gives the following necessary condition. For any
rational prime $q$ splitting in $\mathbb Q(\sqrt{-p})$ and any prime
above it in $K$ of odd residue degree $f$,
\begin{equation}
 4q^f\ge p \quad\text{or}\quad p\mid q^{2f}+q^f+1.
 \tag{CC}\label{eq:cc}
\end{equation}
The second alternative matters for potentially multiplicative reduction.
There is no assumption that $q$ be unramified in $K$.
The same authors, in \href{https://arxiv.org/pdf/2302.08350}
{\emph{Towards strong uniformity for isogenies of prime degree}},
Appendix A, Algorithm A.8 and Example A.13, report the endpoint
$253{,}507$ for the arithmetic test using residue degrees $1,3$. They
apply this test to cubic fields and ask how their method might treat
fields of even degree. Section~\ref{sec:strict-cc} verifies the needed
strict version: every prime $p\equiv3\pmod4$ larger than $253{,}507$
has a prime $q$ splitting in $\mathbb Q(\sqrt{-p})$ with
\[
 4q^3<p,\qquad p\nmid(q^2+q+1)(q^6+q^3+1).
\]
The small split prime bound uses Bach and Sorenson for
$\mathbb Q(\sqrt{-p})$. Its zeta function is
$\zeta(s)L(s,\chi_{-p})$. This explains why Dirichlet GRH
supplies the analytic hypothesis even when $L/\mathbb Q$ is nonabelian.
In using the result one must retain the strict inequality even when
$p$ lies just below $4q^3$. The endpoint $p=3$ has no $q$ with $4q<p$
and so passes the arithmetic test vacuously.

\paragraph{Residue degrees, including ramification.}
Write $G=\operatorname{Gal}(L/\mathbb Q)$, $H=\operatorname{Gal}(L/K)$
and $X=G/H$. At a prime of $L$ over $q$, let $I\triangleleft D$ be the
inertia and decomposition groups. Take a lift $g\in D$ of the Frobenius
generator of $D/I$. The $D$-orbits on $X$ correspond to primes of $K$
over $q$. In each such orbit the inertia orbits are permuted transitively
by $g$, and their number is the residue degree $f$.
Indeed, for a point with stabilizer $D_x$ the number is $[D:I D_x]$,
the residue degree in the associated local subfield. Consequently $f$
divides the length of the $g$-cycle of any point in that orbit.
An odd cycle of length $1$ or $3$ supplies a prime with $f=1$ or $3$.

\begin{center}
\begin{tabular}{lrrrr}
 & $1^6$ & $1^2 2^2$ & $3^2$ & $1\,5$\\
\hline
$A_4$ on edges & 1 & 3 & 8 & 0\\
$A_5$ on Sylow $5$-subgroups & 1 & 15 & 20 & 24
\end{tabular}
\end{center}
Here the entries count group elements of the indicated cycle shape.
For $A_4$ the double transpositions fix two edges. For $A_5$ the
stabilizer of a point is a dihedral group of order $10$. Its five
involutions account, by counting incidences across the six stabilizers,
for two fixed points per involution. Each element of order five fixes
its unique Sylow subgroup. Each element of order three acts without a
fixed point. These observations prove the table and show why the same
arithmetic sieve applies to both sextic actions. Testing the primes
$p\equiv3\pmod4$ up to $253{,}507$ gives $66$ survivors. The complete list
is in \path{checks/isogeny_sieve.out}. This proves the upper bound.

\paragraph{General uniform bound.}
Suppose more generally that every element of a specified transitive
permutation group on $G/H$ has an odd cycle of length at most $m$. Then
each prime of $\mathbb Q$ has an odd residue degree $f\le m$ somewhere
above it in $K$. Write $U(p)=(4\log p+10)^2$. The explicit splitting
bound used in Banwait--Derickx, Appendix A, gives a split prime
$q\le U(p)$ in $\mathbb Q(\sqrt{-p})$. Thus every Momose type 2 prime
must satisfy
\[
 p\le \max\{4U(p)^m,\ U(p)^{2m}+U(p)^m+1\}.
\]
If both expressions were smaller than $p$, both alternatives in
\eqref{eq:cc} would fail for the prime of odd residue degree just found.
The inequality gives an effective bound depending only on $m$, since its
right side is a fixed polynomial in $\log p$.

A variant allows a fixed abelian extension $F/\mathbb Q$ inside $L$.
It suffices that every element of $\operatorname{Gal}(L/F)$ have a short
odd cycle. Apply the least split prime estimate to the abelian field
$J=F\,\mathbb Q(\sqrt{-p})$ instead. Its degree is at most
$2[F:\mathbb Q]$, and its absolute discriminant is at most
$|D_F|^2p^{[F:\mathbb Q]}$ by the compositum discriminant inequality.
The resulting $q\ll_F (1+\log p)^2$ splits in both fields. Its
decomposition group in $L$ acts trivially on $F$, so the same argument
works even if $q$ ramifies in $L/F$. For example, the edge action of
$S_4$ has this property with $m=3$ once its quadratic sign field is fixed.
This yields a bound uniform over these sextics with fixed sign field.

\paragraph{Sharpness of the local obstruction.}
Fix any one of the surviving primes $p$. Let $S$ consist of the primes
$q<p/4$ splitting in $\mathbb Q(\sqrt{-p})$. We first construct a cyclic
cubic field $C$ in which all primes in $S$ are inert. Consider
\[
 u_t(X)=X^3-tX^2-(t+3)X-1,\qquad
 \operatorname{disc}(u_t)=(t^2+3t+9)^2.
\]
For every prime $q$, the values $0,-1$ are never roots. Each of the other
$q-2$ residue classes $x$ gives at most one $t$ with $u_t(x)=0$,
since the coefficient of $t$ is $-x(x+1)$. Thus there is some $t$ with
$u_t$ irreducible modulo $q$. Choose these parameters simultaneously
by the Chinese remainder theorem, adding an auxiliary prime if $S$ is
empty. Then $u_t$ is irreducible over $\mathbb Q$, its splitting field
$C=\mathbb Q(\alpha)$ is cyclic cubic, and every $q\in S$ is unramified
and inert because its reduction is separable and irreducible.

Here are explicit details for extending $C$ to a tetrahedral normal
closure. Choose a rational prime splitting completely in $C$ and use
weak approximation to find $\beta\in C^\times$ with valuations $(1,0,0)$
at its three primes. For a generator $\sigma$ of
$\operatorname{Gal}(C/\mathbb Q)$ set $a=\beta/\sigma\beta$.
The classes of $a$ and $\sigma a$ are independent in
$C^\times/(C^\times)^2$, since their valuation vectors modulo two are
two distinct nonzero even-weight vectors on the three primes.
The relation $a\,\sigma a\,\sigma^2 a=1$ gives a Galois field
$L=C(\sqrt a,\sqrt{\sigma a})$ over $\mathbb Q$. Its order-$3$ quotient
permutes the three nontrivial square classes cyclically. An order-$3$
subgroup exists by Cauchy's theorem, so the group is the nontrivial
semidirect product $V_4\rtimes C_3\cong A_4$. The sextic subfield
$K=C(\sqrt a)$ has the tetrahedral edge action.

For $q\in S$, every residue degree in $K$ is a multiple of $3$ and at
most $6$. Thus the only odd degree to consider is $3$. If $4q^3<p$,
then $q^2+q+1<p$. Since $p$ passed the uniform sieve, we must have
$p\mid q^6+q^3+1$. If $4q^3\ge p$, the first clause in
\eqref{eq:cc} already holds. Every other split prime has $4q\ge p$.
This proves \eqref{eq:cc} for the constructed field at every prime,
establishing the asserted exact classification of compatible primes.

For the endpoint $p=253{,}507$, the set $S$ has $3{,}152$ primes, from
$41$ through $63{,}361$. Already $4\cdot41^3>p$.
For this witness the supplied parameter is rational, $t=A/B$ in lowest
terms. To shorten the defining polynomial, the checker first combines
admissible residues into $t_0$ modulo $M=17^2\prod_{q\in S}q$, with
$t_0\equiv0\pmod{17^2}$. The Euclidean algorithm gives a short pair
$(A,B)$ with $A\equiv t_0 B\pmod M$ and $\gcd(B,M)=1$.
The largest absolute coefficient of the integral polynomial
$B X^3-A X^2-(A+3B)X-B$ has $6{,}796$ decimal digits. A monic integral
polynomial for the same field is
$Y^3-A Y^2-B(A+3B)Y-B^3$, with largest coefficient of $20{,}382$
digits in absolute value. Since the
denominator is a unit at each prescribed prime, reduction and
irreducibility work as before. Modulo $17$, $u_t$ has three simple
roots $3,4,10$ and $u_t(3)\equiv17\pmod{17^2}$. Hence
$\beta=\alpha-3$ has precisely the valuations needed above. One can take
$\sigma(\alpha)=-1/(\alpha+1)$, as substitution into $u_t$ verifies.
The two integers in \path{checks/sextic_parameter.txt} define a specific sextic-field
witness via two short equations. It also gives a cyclic cubic witness
to the local compatibility question discussed in Remark A.14 of the same
preprint. Their cubic example has coefficients exceeding $27{,}380$
decimal digits and enforces inertness at every prime below $p/4$.
The weaker splitting prescriptions used here suffice for the same
necessary condition.

\paragraph{A single explicit tetrahedral field for all 66 primes.}
The preceding construction permits simultaneous prescriptions. Take
$S$ to be the union, over the $66$ candidates $p$, of all split primes
$q<p/4$ in $\mathbb Q(\sqrt{-p})$. The union comprises $5{,}093$ primes
and omits $r=9{,}343$. Require inertness in $C$ at the primes of $S$
and take $t\equiv0\pmod{r^2}$. The polynomial $X^3-3X-1$ has simple
roots $755,8688,9243$ modulo $r$, and its value at $755$ has $r$-adic
valuation one. Thus $\beta=\alpha-755$ gives the required independent
square classes. With the two integers in
\path{checks/simultaneous_parameter.txt}, set
\[
 t=A/B,\quad u_t(\alpha)=0,\qquad
 K=\mathbb Q\!\left(\alpha,
    \sqrt{(\alpha-755)/(-1/(\alpha+1)-755)}\right).
\]
The largest absolute coefficient of $B u_t(X)$ has $10{,}797$ decimal
digits. For this one field, (CC) holds for precisely the $66$ candidates
among primes $p\equiv3\pmod4$. Indeed each relevant $q$ is inert in $C$,
and the finite check gives $4q^3\ge p$ for every such pair $(p,q)$.
The first alternative in (CC) therefore suffices whenever the residue
degree in $K$ is odd. Section~\ref{sec:strict-cc} and the sextic bound
exclude all other primes in the specified congruence class. The verifier
\path{checks/simultaneous_witness.py} rebuilds the union and checks the
parameter at each of its primes, including $2$.

\paragraph{Sharpness also for the primitive sextic fields.}
The same set of 66 primes is exact for local compatibility in the $A_5$
family. Use the same union $S$ of splitting primes over all $66$ candidates. We construct an $A_5$-extension in
which every $q\in S$ is unramified with Frobenius a $3$-cycle in the
natural quintic action. In the sextic action its cycle lengths are $3,3$.
The proof of local compatibility above then applies unchanged.

Let $X$ parametrize monic quintics $f$ together with $y\ne0$ satisfying
$y^2=\operatorname{disc}(f)$. Its function field is
$\mathbb Q(x_1,\ldots,x_5)^{A_5}$: the elementary symmetric functions of
the roots generate the $S_5$-invariant field and the Vandermonde product
generates its quadratic extension. Maeda proved this $A_5$-invariant
field is purely transcendental. Consequently $X$ has weak approximation
at finite sets of primes. Indeed, a nonempty Zariski open of $X$ is
isomorphic to an open in affine $5$-space. Every nonempty local open on
$X(\mathbb Q_q)$ meets this chart since $X$ is smooth and geometrically
integral. Ordinary weak approximation in the five coordinates now
suffices. See \href{https://doi.org/10.1016/0021-8693(89)90174-9}
{T. Maeda, \emph{Noether's problem for $A_5$}, J. Algebra 125 (1989)}.

At every $q\in S$ prescribe squarefree reduction of type $(3,1,1)$.
Such a polynomial exists even over $\mathbb F_2$, by taking an irreducible
cubic and two distinct linear factors. At two further primes prescribe
types $(3,1,1)$ and $(5)$. Each prescription defines a nonempty local
open on $X$. In fact, a monic integral lift of the chosen reduction has
unramified splitting field with Frobenius given by the indicated cycle
type. This permutation is even, so the signed Vandermonde product of its
roots belongs to $\mathbb Q_q$ and supplies $y$. Weak approximation now
gives one rational quintic realizing all prescriptions. Its group is
contained in $A_5$ and has order divisible by $15$. A subgroup of order
$30$ would contradict simplicity. A subgroup of order $15$ would have
normal Sylow $5$-subgroup, whereas the normalizer of this subgroup in
$A_5$ has order $10$. Hence the splitting field has group $A_5$, as required.

\section{Which alternating-group subset fields admit the argument?}
\label{sec:subsets}

Let $L/\mathbb Q$ be a Galois extension with group $A_n$, $n\ge4$. Fix
a $k$-element subset of the $n$ letters, and let $K$ be the fixed field
of its setwise stabilizer. Then $[K:\mathbb Q]=\binom nk$. Complementation
allows us to take $1\le k\le n/2$. The following criterion gives an infinite
sequence of even degrees in which the argument of the preceding section works
uniformly in $L$. For a permutation, orbit length here refers to its action
on the set of all $k$-subsets.

\paragraph{Subset criterion.}
Every element of $A_n$ has an odd orbit if and only if either
$\binom nk$ is odd, or
\[
       n\in\{2^a,2^a+1\},\qquad k=2^{a-1},\qquad a\ge2.
\]
In the displayed cases, the maximum over $g\in A_n$ of the length of its
shortest odd orbit is exactly the largest odd integer $m\le n$.
Consequently every rational prime has an odd residue degree at most $m$
somewhere above it in $K$, including primes ramified in $L$.
The bound of the preceding section applies with $m$, independently of
the discriminant of $K$. If the criterion fails, a positive-density set
of unramified primes has exclusively even residue degrees in $K$.

\paragraph{Proof of the criterion.}
A permutation $g$ has an odd orbit on subsets precisely when the Sylow
$2$-subgroup of $\langle g\rangle$ fixes a subset. This follows by taking
the index of its stabilizer in the cyclic group. It suffices, therefore, to
test all even permutations of $2$-power order. Their cycle lengths form
partitions of $n$ into powers of two, with an even number of even parts.
The fixed subset sizes are exactly the subset sums of the parts.

Write $B$ for the powers that occur in the binary expansion of $n$.
Every partition into powers of two refines $B$, as one sees by merging
equal parts until all parts are distinct. Splitting a part leaves every
previously attainable subset size attainable. Therefore every submask of
$n$ (every sum of a subset of $B$) always works. These sizes have
$\binom nk$ odd, by expanding $(1+X)^n=\prod_{2^j\in B}(1+X^{2^j})$
in $\mathbb F_2[X]$.

Suppose $k$ is not a submask. If $B$ itself describes an even permutation,
it excludes $k$. Otherwise, every even partition is a proper refinement,
and replacing any one part $2^j>1$ of $B$ by two halves already gives an
even permutation. Thus $k$ must be a subset sum after \emph{every}
possible single split. This is sufficient too, because a refinement starts
with one of those splits. For a given $j$, a new sum must use exactly one
of the two halves. If $2^{j-1}\in B$ already, the two copies used in such
a sum either leave the binary bit at $j-1$ or carry into the vacated
bit $j$. All sums are then old. If $2^{j-1}\notin B$, a new sum has
exactly one bit outside $B$, at $j-1$. Two different choices of $j$
therefore cannot allow the same new sum. The only remaining $n$ have
a single binary part greater than one, with an optional part $1$.
For $n\ge4$ the single split allows precisely the additional sizes
$2^{a-1}$ and, when $n$ is odd, its complement $2^{a-1}+1$.
This proves the equivalence. For an excluded pair Chebotarev realizes
the conjugacy class of the excluding element with positive density.

\paragraph{The sharp orbit length.}
Put $k=2^{a-1}$. Write each cycle length of $g$ as $2^bu$ with $u$ odd.
Its $2$-part $2^b$ divides $k$ in both listed degrees. Indeed a cycle
violating this would have length $2^a$ and at most one point outside it,
making the permutation odd. Sort the cycles by decreasing $2^b$. Take
whole cycles until the next cycle would reach or exceed $k$ letters.
The remaining number of letters required is a multiple of the $2^b$ of
this cycle. Its $u$th power splits the cycle into $u$ cycles of length
$2^b$. Take the required number of these. The resulting $k$-set is fixed
by $g^u$, so its orbit has odd length dividing $u\le m$.

When $n=2k$ a cycle of length $m=2k-1$ with a fixed letter attains $m$.
A $k$-set selects $k$ or $k-1$ letters in the long cycle, and both
numbers are coprime to $m$. Its stabilizer in this regular cyclic action
has order dividing the number selected. Hence it is trivial. When
$n=2k+1$ the $n$-cycle attains $m=n$ by the same coprimality argument.
Both witness permutations are even. The inertia-orbit argument in the
preceding section gives the assertion about all rational primes.
\hfill$\square$

\begin{center}
\begin{tabular}{rrrrl}
\toprule
$n$ & $k$ & $[K:\mathbb Q]$ & $m$ & Cycle type attaining $m$\\
\midrule
4 & 2 & 6 & 3 & $3+1$\\
5 & 2 & 10 & 5 & $5$\\
8 & 4 & 70 & 7 & $7+1$\\
9 & 4 & 126 & 9 & $9$\\
16 & 8 & 12,870 & 15 & $15+1$\\
17 & 8 & 24,310 & 17 & $17$\\
\bottomrule
\end{tabular}
\end{center}
For example, a degree-$70$ field in the table can use the same universal
odd-residue-degree test as a degree-seven field. The sharpness just proved
concerns orbit lengths. Determining which primes occur as actual isogeny
degrees requires information about elliptic curves. The script
\texttt{checks/alternating.py} enumerates cycle partitions through $n=34$
and independently visits subset orbits of individual permutations through
$n=7$.

\section{A certificate for the strict local sieve}
\label{sec:strict-cc}
For $p\equiv3\pmod4$ prime, call a rational prime $w$ an excluding witness if
\[
   4w^3<p,\qquad (w/p)=1,\qquad
   p\nmid (w^2+w+1)(w^6+w^3+1).
\]
Quadratic reciprocity identifies $(w/p)=1$ with splitting in
$\mathbb Q(\sqrt{-p})$, including $w=2$. We certify that the entire set
of primes with no witness is as follows.
\[
\begin{gathered}
3,7,11,19,23,31,43,59,67,83,107,139,163,211,283,307,331,379,499,523,\\
547,643,787,883,907,1123,1723,1747,2083,2347,2683,3067,3187,3907,\\
4027,4243,4363,4603,5107,5443,5923,6043,6427,6763,7723,8563,8803,\\
9067,10627,10987,13267,14107,14347,14683,15667,16963,22003,26203,\\
32323,34483,36523,37123,89083,90787,222643,253507 .
\end{gathered}
\]
Trial of the eligible auxiliary primes checks this list through $253507$.
\paragraph{Independent check of the endpoint in the cubic sieve.}
A rejecting auxiliary prime must satisfy $4q^3<p$.
The loop in Banwait--Derickx, Algorithm~A.8, also permits
$q\le 1+(p/4)^{1/3}$, which can include an extra prime.
The authors' historical
\href{https://github.com/isogeny-primes/isogeny-primes/blob/17bb625acaa7bf79dd047a2f01ada1ae725cc582/gp_scripts/partype2primesuniform.gp}{PARI/GP implementation}
uses the same upper endpoint. For example, $p=89083$ passes the strict
criterion. The extra $q=29$ has $(29/p)=1$ and neither polynomial vanishes
modulo $p$, yet $4\cdot29^3=97556>p$. We therefore ran an independent sieve with
the strict integer inequality. For every prime $253507<p\le 11\cdot10^{12}$,
$p\equiv3\pmod4$, the program finds, or certifies the existence of, a prime
$q$ with
\[
 4q^3<p,\qquad (q/p)=1,\qquad
 p\nmid(q^2+q+1)(q^6+q^3+1).
\]
It checks 264911 primes directly up to $4\cdot127^3=8193532$.
Above this threshold all primes $q\le127$ have admissible size. The program
separately checks the eight primes $p$ in its upper range dividing one of
the two polynomials for some $q\le127$ and congruent to $3\pmod4$.
For all remaining primes, a square among these auxiliaries is a witness.
If every auxiliary is a nonsquare, the conditions at $2,3,5,7,11,13,17$
restrict $p$ to 1440 residue classes modulo $2042040$. For every odd $q$,
quadratic reciprocity gives $(q/p)=(-p/q)$. In each class, the further
tests for $19\le q\le127$ are consequently periodic masks on an arithmetic
progression. A mask discards multiples of $q$ and integers with
$(-p/q)=1$. Exactly 327 integers in the upper range survive all the masks.
Trial division proves that 110 are prime, and all 110 have a strict witness.
The complete C++ source and its counts are supplied as
\texttt{checks/cubic\_tail.cpp} and \texttt{checks/cubic\_tail.out}.


\paragraph{An integer certificate for all large primes.}
Let $T=10^{11}$, $y=\log p$ and $M(y)=1.881y+6.18$.
For $p>T>e^{25}$, the quadratic split-prime bound gives a prime
$q\le M(y)^2$ that splits in $\mathbb Q(\sqrt{-p})$. See
\href{https://arxiv.org/pdf/1803.02056}{Elsenhans--Kl\"uners--Nicolae,
Theorem 3}, using Bach--Sorenson and the two Dirichlet factors.
Here $4q^3<p$. Indeed, $M(\log T)<54$ and $4\cdot54^6<T$;
$M(y)^6e^{-y}$ decreases for $y\ge\log T$ by logarithmic differentiation.
If $y\ge58$, the same argument applied to exponents $12$ and $6$, using
$M(58)<116$ and $116^{12}+116^6+1<e^{58}$, yields $q^6+q^3+1<p$.
This $q$ is already a witness.

In the remaining case $q\le116^2$ and $q^2+q+1<T$. The only obstruction
to using $q$ is $p\mid N_q=q^6+q^3+1$. For each of the $1594$ primes
$q\le116^2$ the supplied certificate decomposes $N_q=\prod_i d_i^{e_i}$.
A prime $p>T$ dividing $N_q$ divides some $d_i>T$. For every such entry
the certificate supplies a prime $w\le43$ and an integer $z$ such that
\[
 z^2\equiv w\pmod {d_i},\qquad
 \gcd\bigl(d_i,w(w^2+w+1)(w^6+w^3+1)\bigr)=1 .
\]
Reducing modulo $p$ gives a nonzero square $w$ and both required
nondivisibilities. Since $4w^3<T$, this $w$ is an excluding witness.
This argument works even when $d_i$ is composite. The independent
verifier checks every product, square, gcd and endpoint inequality with
integer or rational arithmetic. It covers 1126 distinct large divisors.
No primality claim about the large divisors enters the proof.
Together the two overlapping ranges prove the classification.

\section{How much splitting information can congruences recover?}
\label{sec:rigidity}

A common route from Dirichlet GRH to arithmetic over number fields is to
express a Dedekind zeta function as a product of Dirichlet $L$-functions.
There is an exact obstruction to applying this route to a nonabelian normal
closure. It can be measured either by the frequency of incorrect splitting
predictions or by the mean squared error in the first Euler coefficients.
The formulas below follow from Chebotarev and Kronecker--Weber. We claim
no priority for them. For the classical splitting problem, see
\href{https://arxiv.org/abs/2208.11381}{Rosen, Takeyama, Tasaka and Yamamoto};
for the input theorems see
\href{https://www.jmilne.org/math/CourseNotes/cft.html}{Milne, \emph{Class Field Theory}}.

\paragraph{Best possible congruence test.}
Let $L/\mathbb Q$ be finite Galois with group $G$, let $D=[G,G]$, and let
$C\subseteq G$ be a union of conjugacy classes. One wants to decide whether
$\operatorname{Frob}_p\in C$ from a fixed residue $p\bmod m$. Allow
\emph{any} fixed modulus $m$ and \emph{any} choice of acceptable residue
classes. The minimum natural density, relative to all primes, of incorrect
decisions is exactly
\[
 \boxed{\displaystyle
  \frac{1}{|G|}\sum_{T\in G/D}\min\{\,|T\cap C|,\ |T\setminus C|\,\}.}
\]
The minimum is attained. Ramified primes and primes dividing the modulus
form a finite exceptional set.

\emph{Proof.} Put $M=\mathbb Q(\zeta_m)$ and $E=L\cap M$.
The Galois group of $LM$ consists of the compatible pairs in
$G\times\operatorname{Gal}(M/\mathbb Q)$ with equal restriction to $E$.
By Chebotarev, prime densities for any conjugacy-invariant test on these
pairs are uniform element counts. Fixing the second coordinate leaves the
first uniformly distributed, for such tests, in a coset of
$\operatorname{Gal}(L/E)$. Since $E$ is abelian, this coset is a union
of $D$-cosets. Even if one learns the individual $D$-coset, the majority
decision on it still makes at least the displayed number of errors.
Kronecker--Weber supplies an $m$ for which $L^D\subseteq M$. Then
$E=L^D$ and the residue reveals exactly the required coset. Choosing its
majority value attains the bound.\hfill$\square$

For example, let $K$ be any number field and take $L$ to be its normal closure.
A prime splits completely in $K$ precisely when its Frobenius acts trivially
on all embeddings of $K$. This action is faithful, so $C=\{1\}$. The minimum
error density equals $0$ for abelian $G$ and $1/|G|$ for nonabelian $G$.
In the second case the rule declaring that no prime splits completely attains
the minimum. In particular, even the best rule for complete splitting in
$\mathbb Q(\sqrt[3]{2})$ makes errors on a set of density at least $1/6$.
The splitting field has group $S_3$, since $X^3-2$ is irreducible by Eisenstein
and has nonsquare discriminant $-108$.

\paragraph{Best approximation to the Euler coefficient.}
Write $H=\operatorname{Gal}(L/K)$ and $F=L^{HD}$, the largest subfield of $K$
abelian over $\mathbb Q$. For an unramified prime $p$, write $a_K(p)$ for the
number of degree-one primes of $K$ above $p$. The optimum over all
complex-valued periodic functions $b$ is
\[
 \boxed{\displaystyle
 \min_b\lim_{x\to\infty}\frac{1}{\pi(x)}\sum_{p\le x}
          |a_K(p)-b(p)|^2
      = |H\backslash G/H|-[F:\mathbb Q].}
\]
An optimal $b$ equals $a_F(p)$ away from the finite exceptional set.

\emph{Proof.} Let $u(g)$ count the fixed points of $g$ on $G/H$. At every
unramified $p$ this is $a_K(p)$ for $g=\operatorname{Frob}_p$. By the
same compatible-pair argument, the sharpest available conditional mean
is the average of $u$ on each $D$-coset, attained at a modulus containing
$L^D$ in its cyclotomic field. Indeed a squared error on a finite set is
minimized by its mean, and a coarser partition cannot reduce the minimum.
The coset average is the character of the $D$-invariants of the permutation
representation on $G/H$, because averaging the $D$-action projects onto
this subspace. A basis is indexed by $D$-orbits, which form the
$G$-set $G/HD$. Its character is precisely $a_F$ on Frobenii. This
character is $[G:HD]$ on $HD$ and zero elsewhere, since $HD$ is normal.
Its squared norm for normalized counting measure on $G$ equals $[G:HD]$.
The squared norm of $u$ counts $G$-orbits on $(G/H)^2$ by Burnside's lemma.
Those orbits correspond to the $H$-double cosets in $G$. Subtracting
the norm of the orthogonal projection gives the formula.\hfill$\square$

For $\mathbb Q(\sqrt[3]{2})$ the best approximation is the constant $1$.
The values $3,1,0$ occur with densities $1/6,1/2,1/3$, giving squared
error $4/6+1/3=1$. The sign of Frobenius, which gives the quadratic
resolvent, distinguishes odd from even permutations but does not improve
this squared error. The file \texttt{checks/rigidity.py} checks these
finite group calculations and the corresponding ones for cyclic quartic,
dihedral quartic and $S_4$ quartic permutation actions.

\paragraph{Consequence for direct transfer.}
Suppose $\zeta_K$ agrees in $\Re s>1$ with a finite product of integer powers
of ordinary Dirichlet $L$-functions up to Euler factors at finitely many primes.
By uniqueness of Dirichlet coefficients, its coefficient at almost every $p$
equals a fixed linear combination of values $\chi(p)$. This is periodic in $p$,
so the condition $a_K(p)=[K:\mathbb Q]$ would predict complete splitting
with zero error density. The first formula forces $G$ to be abelian.
Conversely, abelian fields over $\mathbb Q$ have precisely such a
factorization via their primitive Dirichlet characters. Thus any attempt
to establish GRH for the zeta function of a nonabelian normal closure,
or of its nonnormal generating subfield, needs an argument beyond a
finite Dirichlet product or quotient. This obstruction concerns that
particular proof strategy. Zeros of distinct analytic functions can still
be related by other means.

\paragraph{A one-sided criterion for the missing zeros.}
There is a concrete use for the optimal predictor $a_F$. Put
\[
 U_K(x)=\sum_{p\le x,\ p\text{ unramified in }L}
       (a_K(p)-a_F(p))\log p\,(1-p/x).
\]
With the repository's Dirichlet GRH as input, the following are equivalent:
\[
 \mathrm{RH}(\zeta_K)
 \quad\Longleftrightarrow\quad
 \exists C_K\ \ \forall n\in\mathbb Z,\ n\ge2:\quad
 U_K(n)\ge-C_K\sqrt n .
\]
This works for nonnormal $K$ as well. In the primitive $A_5$ sextic
case, $F=\mathbb Q$ since $A_5$ is perfect. The weights $a_K(p)-1$
are respectively $5,1,-1,0$ at Frobenius elements of orders
$1,2,3,5$. Thus the missing inequality is a single lower bound for a
sum over rational primes with four explicit weights.

\emph{Proof.} Write $V_J(x)=\sum_{N\mathfrak p^r\le x}\log N\mathfrak p
(1-N\mathfrak p^r/x)$, summing over prime ideals of a field $J$ and
integers $r\ge1$. Then $V_J=x^{-1}\int_1^x\psi_J(u)\,du$. Under field RH
the integrated explicit formula yields $V_J=x/2+O_J(\sqrt x)$: the
nontrivial zero sum after division by $x$ is bounded by
$\sqrt{x}\sum_\rho |\rho(\rho+1)|^{-1}$, a convergent sum.
See \href{https://arxiv.org/abs/1312.4465}{Greni\'e--Molteni,
formula (1.3a)}. Prime powers of rational exponent at least two
make a contribution $O_J(\sqrt x)$, by bounding their log coefficients
by $[J:\mathbb Q]\log p$ and using the ordinary Chebyshev bound,
first for squares and then for exponents at least three. The finitely many
ramified primes cost $O_J(\log x)$. Thus $U_K=V_K-V_F+O_L(\sqrt x)$.
Here $F$ satisfies RH by abelian factorization. This proves necessity.

For sufficiency extend the integer lower bound to all real $x\ge2$.
Between primes $|U_K'(x)|\ll_K x^{-2}\sum_{p\le x}p\log p\ll_K 1$,
and $U_K$ is continuous. Write $W=V_K-V_F$. For $\Re s>1$,
\[
 \int_1^\infty W(x)x^{-s-1}\,dx
 =\frac{-\zeta'_K(s)/\zeta_K(s)+\zeta'_F(s)/\zeta_F(s)}{s(s+1)} .
\]
The pole terms at $1$ cancel. Suppose $\zeta_K$ has a zero $\rho$ with
$\Re\rho>1/2$. It gives a genuine pole since $\zeta_F$ is zero-free there.
Add $C\sqrt x$ to $W$ to make it nonnegative, increasing $C$ on a compact
interval if necessary. This adds $C/(s-1/2)$ to the transform and preserves
the pole at $\rho$. The integral's real abscissa of convergence $b$
therefore satisfies $\Re\rho\le b\le1$. Landau's theorem forces a
singularity at the real point $b$: Taylor expansion from the right has
nonnegative moments of $\log x$ as coefficients in the direction of
decreasing $s$, so analyticity at $b$ would give convergence across $b$.
At $b=1$ the transform is regular. At any real zero $b\in(1/2,1)$ of
multiplicity $h$, its residue is $-h/[b(b+1)]<0$, inconsistent with the
nonnegative integral to the right. These exhaust the possible real
singularities. The functional equation gives the other half of the
critical strip.\hfill$\square$

The incoming reports apply a related sign argument to completely split
primes in a Galois field and to individual cubic examples. The conditional
mean here selects the abelian contribution for every field and the smoothing
removes the pole without a field-dependent centering constant. Proving the
displayed bound for the primitive sextic family still requires control of
the nonabelian weights.


\section{Further research questions}
\begin{enumerate}
\item Is there a prime $q$ for which the first prime $1\bmod q$ is already
$1\bmod q^2$? The geometric-interval argument makes this precise question
equivalent to Pomerance's divisibility conjecture failing. The first interval
of Barthel--M\"uller's conjecture would rule out every such $q$.
\item Which of the $66$ local candidates actually occur as degrees of
elliptic-curve isogenies over the specified sextic fields? The constructions
show that the local condition, even when tested at every rational prime,
cannot further shorten the list uniformly. The next obstruction must use
additional arithmetic of the curves or of their moduli.
\item For which further transitive Galois actions do suitable abelian
splitting conditions force short odd Frobenius orbits? A positive answer
gives more even-degree uniformity bounds. The subset theorem gives an exact odd-orbit criterion for the alternating-group actions considered here.
\item Can a theta-lift or height estimate bound $U_K(x)$ from below in
the primitive $A_5$ sextic setting? The explicit weights and the best
periodic approximation quantify the missing arithmetic information.
\end{enumerate}
\end{document}
